\BeleskeID{09_3-s011}
% element: 09_3-s011:e04
% slide-id: 09_3-s011
Pri jednakim početnim širinama $W_{n2}=W_{p2}$ uzima se $C_{Gp2(\beta=1)}=C_{Gn2}$. Povećanjem širine P kanala njegova gejtska kapacitivnost raste približno linearno, dok se ekvivalentna otpornost smanjuje približno obrnuto proporcionalno širini.

Srednje kašnjenje dobijamo sabiranjem doprinosa za dva smera prelaza:
\[\begin{aligned}
t_p&=\frac{0.69}2(R_{n1}+R_{p1})\bigl((C_{Dp1}+C_{Dn1})+(C_{Gp2}+C_{Gn2})\bigr)\\
&=\frac{0.69}2\left(R_{n1}+\frac{R_{p1(\beta=1)}}\beta\right)
\bigl((\beta C_{Dp1(\beta=1)}+C_{Dn1})+(\beta C_{Gp2(\beta=1)}+C_{Gn2})\bigr)\\
&=\frac{0.69}2R_{n1}\left(1+\frac r\beta\right)(1+\beta)(C_{Dn1}+C_{Gn2}).
\end{aligned}\]
U korišćenom izrazu za otpornosti odnos je
\[r=\frac{\frac34\frac{V_{DD}}{I_{Dpsat}}(1-\frac79|\lambda_p|V_{DD})}
{\frac34\frac{V_{DD}}{I_{Dnsat}}(1-\frac79\lambda_nV_{DD})}.
\]
Pri posmatranju promene $\beta$ ostali parametri u ovom izvođenju uzimaju se kao fiksni.
